\documentclass{leveldoc}
\DocUses{foundations,first_bend,slingshot,around_the_wall,einzel_lens,injection}
\DocFingerprint{01893770e72e06e9}

\begin{document}

\section{What the level asks}\label{sec:asks}

The game's text:
\leveldescription

Three shots leave $(\val{xLaunch}, \val{yLaunch})$ with the same particle
($q = \val{q}$, $m = \val{m}$) and the same kinetic energy $T_0 = \val{Tzero}$.
The directions stored in the file are unit vectors at exactly $\val{angDown}^\circ$,
$\val{angAxis}^\circ$ and $\val{angUp}^\circ$. On every shot the detector is the
rectangle $\val{xDlo} \le x \le \val{xDhi}$, $\val{yDlo} \le y \le \val{yDhi}$.
The box is tall. The requirement is the direction: an entry counts only when the
velocity lies inside the cone of \xeqref{injection}{eq:cone}, of half-angle
$\val{halfDeg}^\circ$ about the file's axis at $\val{axisDeg}^\circ$. The player
places at most $\val{maxCharges}$ charges, of the particle's sign only, inside the
region $\val{xRegLo} \le x \le \val{xRegHi}$, $\val{yRegLo} \le y \le \val{yRegHi}$,
and no magnets. A flight ends when the particle touches a charge, leaves the world,
enters the detector or runs out of time
(\xsecref{foundations}{sec:board}); $t_{\max} = \val{tmax}$.

The reference solution adds $Q = \val{Qlo}$ at $(\val{xLo}, \val{yLo})$ and
$Q = \val{Qhi}$ at $(\val{xHi}, \val{yHi})$. Both have the particle's sign. The lower
charge is the stronger one and sits below the beam; the upper charge sits above it.
Together they bend the three rays back toward $+x$, and each arrives inside the cone.
This document integrates those two-charge flights, and the flight of each charge
alone, by the fixed-step method of \levelref{around_the_wall}, and applies the cone
at the entry.

\section{Setup}\label{sec:setup}

The level file, field by field (the appendix reproduces it):
\setuptable

In symbols: $c = \val{c}$, the particle has radius $a = \val{aRad}$, and every charge
has radius $R = \val{R}$. The launch has
\[
  \gamma_0 = 1 + \frac{T_0}{mc^2} = \val{gammaZero}, \qquad
  \abs{\vect{p}_0} = \sqrt{\frac{T_0^2}{c^2} + 2mT_0} = \val{pZero}, \qquad
  v_0 = \frac{\abs{\vect{p}_0}}{\gamma_0 m} = \val{vZero}.
\]
The same $\gamma_0$, $\abs{\vect{p}_0}$ and $v_0$ start every shot; only the direction
changes. The world is $-1 \le x \le \val{xWorld}$, $-1 \le y \le \val{yWorld}$.
The couplings are $\kappa_\mathrm{l} = qQ_\mathrm{l} = \val{kLo}$ and
$\kappa_\mathrm{h} = qQ_\mathrm{h} = \val{kHi}$, in the file's order. The cone's axis
is $+x$ and its half-angle in radians is $\delta = \val{halfRad}$.

\begin{figure}[htbp]
\centering
\begin{tikzpicture}
\begin{axis}[axis equal image, width=0.88\linewidth, xmin=-1.8, xmax=32.4, ymin=-1.8, ymax=21.8,
  xtick={0,5,...,30}, ytick={0,5,...,20}, tick label style={font=\small},
  xlabel={$x$}, ylabel={$y$}, label style={font=\small}]
  \draw[black!45, dashed] (-1,-1) rectangle (31,21);
  \draw[black!50] (0,0) rectangle (30,20);
  \draw[black!40, dashed] (4,1) rectangle (20,19);
  \fill[green!50!black, fill opacity=0.16] (28,6) rectangle (30,14);
  \draw[green!40!black] (28,6) rectangle (30,14);
  \fill[red!70!black, fill opacity=0.35] (\rawval{xLo},\rawval{yLo}) circle[radius=0.3];
  \draw[red!70!black] (\rawval{xLo},\rawval{yLo}) circle[radius=0.3];
  \fill[red!70!black, fill opacity=0.35] (\rawval{xHi},\rawval{yHi}) circle[radius=0.3];
  \draw[red!70!black] (\rawval{xHi},\rawval{yHi}) circle[radius=0.3];
  \addplot[thick, orange!80!black] table {pathdown.dat};
  \addplot[very thick, blue!70!black] table {pathaxis.dat};
  \addplot[thick, purple!70!black] table {pathup.dat};
  \draw[-{Stealth[length=6pt]}, thick, orange!80!black] (\rawval{xLaunch},\rawval{yLaunch}) --
    (\rawval{arrowDx},\rawval{arrowDy});
  \draw[-{Stealth[length=6pt]}, very thick, blue!70!black] (\rawval{xLaunch},\rawval{yLaunch}) --
    (\rawval{arrowX},\rawval{arrowY});
  \draw[-{Stealth[length=6pt]}, thick, purple!70!black] (\rawval{xLaunch},\rawval{yLaunch}) --
    (\rawval{arrowUx},\rawval{arrowUy});
  \fill (\rawval{xLaunch},\rawval{yLaunch}) circle[radius=0.11];
  \node[font=\scriptsize, red!60!black, anchor=north] at (4,2.55) {$Q_\mathrm{l}$};
  \node[font=\scriptsize, red!60!black, anchor=south] at (8,16.45) {$Q_\mathrm{h}$};
  \node[font=\scriptsize, green!30!black, anchor=west] at (30.15,10) {detector};
\end{axis}
\end{tikzpicture}
\caption{The three rays of the reference. Orange is $\val{angDown}^\circ$, blue is the
axis, purple is $\val{angUp}^\circ$. All three enter the left face of the detector
inside the cone. The discs are the two repelling charges, each of radius $R$. The
dashed rectangle is the region in which the player may place charges.}
\label{fig:board}
\end{figure}

\section{Principles and model}\label{sec:model}

\begin{principles}
  \item The field of a charged sphere and the force on a non-polarizable particle,
    \xeqref{foundations}{eq:coulomb} and \xeqref{foundations}{eq:sphereforce}.
  \item The relativistic equation of motion \xeqref{foundations}{eq:eom}, the launch
    \xeqref{foundations}{eq:launch} and the energy \xeqref{foundations}{eq:energy}. The
    motion stays in the plane (\xsecref{foundations}{sec:plane}).
  \item The board's rules (\xsecref{foundations}{sec:board}). The first entry into the
    detector decides, and that entry has to lie in the cone of
    \xeqref{injection}{eq:cone}. This level adds no kinetic window.
  \item The classical Runge--Kutta method \xeqref{foundations}{eq:rk4} and the fixed-step
    flight of \xsecref{around_the_wall}{sec:method}. With two centres the orbit equation
    does not separate (\xeqref{slingshot}{eq:noseparation}).
  \item The radiated power \xeqref{foundations}{eq:larmor}, with the acceleration
    \xeqref{first_bend}{eq:acceleration}, summed by Simpson's rule
    \xeqref{foundations}{eq:simpson}.
\end{principles}

The level has no magnets, so $\vect{B} = \vect{0}$. The potential and the kinetic energy
are \xeqref{einzel_lens}{eq:potential} and \xeqref{einzel_lens}{eq:kinetic}, here with
$\kappa_\mathrm{l} = \val{kLo}$ at $(\val{xLo}, \val{yLo})$ and
$\kappa_\mathrm{h} = \val{kHi}$ at $(\val{xHi}, \val{yHi})$. The force is
\xeqref{around_the_wall}{eq:force} with $N = \val{nCharges}$, the charges in the file's
order. The state $(x, y, p_x, p_y)$ obeys \xeqref{foundations}{eq:eom} with that force
until the first event.

\begin{assumption}{The charges are held in place}\label{as:fixed}
None of the charges recoils. That is a rule of the game
(\xsecref{foundations}{sec:radiation}).
\end{assumption}

\begin{assumption}{No radiation}\label{as:radiation}
The radiated energy is not taken from the particle. Over the nine flights of this
document the largest fraction of $T_0$ is $\val{radFracMax}$ (§\ref{sec:hand}). The
game makes the same approximation and reports the same integral (\physref{8}).
\end{assumption}

\begin{assumption}{A point particle}\label{as:point}
$a = 0$, and the particle does not polarize, so \xeqref{foundations}{eq:sphereforce} is
exact. On the three reference rays the closest approach to either charge is
$\val{rClose}$, on the \val{rCloseRay} ray, to the \val{rCloseCharge} charge. The
clearance outside the sphere is $\val{clearClose}$. The samples that locate it are at
most $\val{sampleArc}$ apart (§\ref{sec:method}), far below that clearance.
\end{assumption}

\begin{assumption}{The step resolves the bend}\label{as:step}
The reported step is $h = \val{H}$. Halving it estimates the largest error of an end
state as $\val{errMax}$ (§\ref{sec:hand}). That is far below the half-angle of the
cone, $\val{halfDeg}^\circ$. The tolerance used in §\ref{sec:comparison} is the larger
of $10^{-9}$ and ten times this estimate.
\end{assumption}

\section{Derivation}\label{sec:derivation}

\subsection{What the energy allows}

At the launch, \xeqref{einzel_lens}{eq:potential} gives $V(\vect{x}_0) = \val{VZero}$,
so
\[
  W = \val{W}, \qquad W - mc^2 = T_0 + V(\vect{x}_0) = \val{Wbind}.
\]
This is positive. Far from both charges $V$ falls to $0$, which lies below the bound,
so \xeqref{einzel_lens}{eq:kinetic} allows the particle to leave the world. The
integration has to decide whether it does. On the axis ray the force at the launch is
$(\val{FxZero}, \val{FyZero})$. The horizontal part is negative: both charges stand
ahead of the particle and repel, so the pair retards the beam. The lower charge sits
below the launch and pushes upward; the upper charge sits above it and pushes downward.
The lower coupling is the larger one, and the sum $\val{FyZero}$ points upward.

Both couplings are positive, so $V$ is positive everywhere and largest near the charges.
Along a ray that passes them, \xeqref{einzel_lens}{eq:kinetic} falls and then rises
again. On the axis ray the minimum, computed in §\ref{sec:hand}, is $\val{TminAxis}$,
below $T_0$. At the detector the potential has fallen below its value at the launch,
and the arrival kinetic energies lie above $T_0$. The horizontal momentum stays
positive: the rays are bent, and they do not turn back.

\subsection{The cone}\label{sec:cone}

The acceptance is \xeqref{injection}{eq:cone}. The file's axis is $\val{axisDeg}^\circ$,
so $\uvec{e} = (1, 0)$, and the half-angle is $\delta = \val{halfRad}$. With
$p_x > 0$ the angle from the axis is $\beta = \operatorname{atan2}(\abs{p_y}, p_x)$,
in radians. The entry counts when $\beta \le \delta$. When the particle has entered the
box and $\beta > \delta$, the outcome is Rejected and the reported state stays the
entry (\physref{6.1}). A flight that never enters the box is not tested.

On the three reference rays $p_y > 0$ at the entry, so $\beta$ is the arrival angle
itself. The angles are $\val{betaDown}^\circ$, $\val{betaAxis}^\circ$ and
$\val{betaUp}^\circ$, and the margins inside the half-angle are $\val{marginDown}^\circ$,
$\val{marginAxis}^\circ$ and $\val{marginUp}^\circ$. All three margins are positive.
The axis ray is the closest to the edge of the cone.

\subsection{Why the flight is integrated}\label{sec:noseparation}

With one charge, angular momentum closes the motion into the orbit of
\levelref{first_bend}. Two fixed centres already fail to separate in the relativistic
problem: the term \xeqref{slingshot}{eq:noseparation} does not split. The pair here is
staggered and unequal, so it has no symmetry that would remove the obstruction. Each
ray is the initial-value problem for \xeqref{around_the_wall}{eq:force} at
$N = \val{nCharges}$, and §\ref{sec:method} computes it to a stated accuracy.

\subsection{The method}\label{sec:method}

The state is advanced by the classical Runge--Kutta method \xeqref{foundations}{eq:rk4}
at the fixed step $h$, and a boundary is crossed on the cubic Hermite interpolant
\xeqref{slingshot}{eq:hermite}, as specified in \xsecref{around_the_wall}{sec:method}.
After each step the interpolant is sampled at $\val{nSample}$ points. A sphere, the
world's edge or the detector is tested only when the box of the step's ends, enlarged
by $ch$, meets it. The samples are read in the board's order
(\xsecref{foundations}{sec:board}). The first instant of the winning event is the
bisection of the interpolant, $\val{nBisect}$ iterations. The position reported is the
interpolant at that instant, and the face coordinate is then set to the face value.
The momentum is one further step of the method, of that length. The cone of
§\ref{sec:cone} is applied to that momentum. The error of the reported step is a
fifteenth of its difference from the step $2h$
(\xsecref{foundations}{sec:numerics}).

On each step the same interpolant is also sampled at nine equally spaced fractions.
\xeqref{einzel_lens}{eq:kinetic} at those samples gives the kinetic-energy dip, and
the distances to the two centres give the closest approach. The samples are an eighth
of a step apart in time, so a point between them lies within $\val{sampleArc}$ of a
sample. Those two quantities are located to that tolerance. The end state is held to
the step-halving estimate.

\section{The reference solution by hand}\label{sec:hand}

\begin{procedure}
  \item\label{st:launch} \emph{The launch} (\xeqref{foundations}{eq:launch}):
    $\gamma_0 = \val{gammaZero}$ and $\abs{\vect{p}_0} = \val{pZero}$ on every shot,
    aimed at $\val{angDown}^\circ$, along $+x$, and at $\val{angUp}^\circ$. On the axis
    ray the force is $(\val{FxZero}, \val{FyZero})$, retarding and upward.
  \item\label{st:energy} \emph{The energy}: $W = \val{W}$ and
    $W - mc^2 = \val{Wbind} > 0$. The world's edge is energetically open.
  \item\label{st:flight} \emph{The axis flight}: the method of §\ref{sec:method} with
    $h = \val{H}$. The states every five time units, and the drift of the energy:
    \begin{center}\small\val{states}\end{center}
    The kinetic energy falls to $\val{TminAxis}$ at $t = \val{tTmin}$, at
    $(\val{xTmin}, \val{yTmin})$. That point is $\val{rDipLo}$ from the lower charge
    and $\val{rDipHi}$ from the upper charge. The lowest kinetic energies on the outer
    rays are $\val{RefDownTmin}$ and $\val{RefUpTmin}$.
  \item\label{st:entry} \emph{The three entries.} The axis step from
    $t_n = \val{tBefore}$, at $(\val{xBefore}, \val{yBefore})$, ends at
    $(\val{xAfter}, \val{yAfter})$. The interpolant crosses the \val{faceE} face at the
    fraction $u = \val{uE}$. The residual before the face coordinate is assigned is
    $\val{faceRes}$, so $x_\mathrm{e} = \val{xE}$, $y_\mathrm{e} = \val{yE}$ and
    $t_\mathrm{e} = \val{tE}$, after $\val{steps}$ steps. One step of length
    $t_\mathrm{e} - t_n$ gives $\vect{p}_\mathrm{e} = (\val{pxE}, \val{pyE})$, at
    $\val{angE}^\circ$ to $+x$.

    The $\val{angDown}^\circ$ ray enters the \val{RefDownFace} face at
    $(\val{RefDownX}, \val{RefDownY})$ at $t = \val{RefDownT}$, with momentum
    $(\val{RefDownPx}, \val{RefDownPy})$, at $\val{RefDownAng}^\circ$.

    The $\val{angUp}^\circ$ ray enters the \val{RefUpFace} face at
    $(\val{RefUpX}, \val{RefUpY})$ at $t = \val{RefUpT}$, with momentum
    $(\val{RefUpPx}, \val{RefUpPy})$, at $\val{RefUpAng}^\circ$.

    All three ordinates lie between $\val{yDlo}$ and $\val{yDhi}$, and every
    horizontal momentum is positive.
  \item\label{st:cone} \emph{The cone} \xeqref{injection}{eq:cone}. The three arrival
    angles are the angles $\beta$ of §\ref{sec:cone}, and the margins are
    $\val{marginDown}^\circ$, $\val{marginAxis}^\circ$ and $\val{marginUp}^\circ$.
    Each is positive, so each entry counts.
  \item\label{st:check} \emph{The checks}: along the axis flight the energy stays within
    $\val{wDev}$ of $W$. The arrival kinetic energies from
    \xeqref{einzel_lens}{eq:kinetic} are $\val{TarrDown}$, $\val{TarrAxis}$ and
    $\val{TarrUp}$, each above $\val{Tzero}$. Halving the step estimates the three end
    states as $\val{RefDownErr}$, $\val{errAxis}$ and $\val{RefUpErr}$. The same
    outcomes appear at $h$ and at $2h$.
  \item\label{st:rad} \emph{The radiation}: the power \xeqref{foundations}{eq:larmor}
    with the acceleration \xeqref{first_bend}{eq:acceleration}, summed by Simpson's rule
    on the steps of $h/2$. On the axis ray the integral is $\val{Erad}$, the fraction
    $\val{radFrac}$ of $T_0$. The relative difference from the same sum at step $h$,
    divided by fifteen, is $\val{radErr}$. The outer rays radiate $\val{RefDownRad}$
    and $\val{RefUpRad}$.
\end{procedure}

\begin{result}[A beam inside the cone]
All three rays enter the detector through its left face, travelling within
$\val{halfDeg}^\circ$ of $+x$. The $\val{angDown}^\circ$ ray arrives at
$(\val{RefDownX}, \val{RefDownY})$ at $t = \val{RefDownT}$, at $\val{RefDownAng}^\circ$.
The axis ray arrives at $(\val{xE}, \val{yE})$ at $t = \val{tE}$, at $\val{angE}^\circ$.
The $\val{angUp}^\circ$ ray arrives at $(\val{RefUpX}, \val{RefUpY})$ at
$t = \val{RefUpT}$, at $\val{RefUpAng}^\circ$. The level is solved.
\end{result}

By hand, Step~\ref{st:flight} is $\val{steps}$ steps of four evaluations of the force,
for the axis ray, and the same method for the two outer rays. Across all nine flights
of this document the largest step-halving estimate is $\val{errMax}$.

\section{Comparison with the game}\label{sec:comparison}

\comparisontable

The nine outcomes agree, and so do the nine end positions and times. The reference
momenta are compared as well. The tolerance on a state is $\val{tolFlight}$. Ten times
the largest step-halving estimate $\val{errMax}$ lies below $10^{-9}$, so the tolerance
is that floor. The three reference radiated energies are compared relatively, with
tolerance $10^{-4}$. The state tolerance and the radiation tolerance were fixed from
these estimates before the comparison.

The largest absolute difference is the time of the lower charge's $\val{angUp}^\circ$
ray, $\val{cmpLowUpTDiff}$. Its abscissa differs by $\val{cmpLowUpXDiff}$. The
difference is many times $\val{errMax}$ and a small fraction of the tolerance. The
game's verified end states are good to about $10^{-10}$
(\xsecref{foundations}{sec:verification}), so a difference of this size is what that
accuracy allows when two centres are summed.

The three reference abscissae agree exactly, as do the lower charge's
$\val{angDown}^\circ$ and axis abscissae, the lower charge's $\val{angUp}^\circ$
ordinate, and the upper charge's axis and $\val{angUp}^\circ$ abscissae. The upper
charge's $\val{angDown}^\circ$ abscissa differs by $\val{cmpHighDownXDiff}$.

The axis ray's quadrature uncertainty is $\val{radErr}$, and the reference outer rays
give $\val{RefDownRadErr}$ and $\val{RefUpRadErr}$, all inside the relative tolerance.
The largest relative difference is the reference $\val{angDown}^\circ$ ray,
$\val{cmpRefDownRadDiff}$.

\section{Other solutions}\label{sec:others}

Either added charge, without the other, fails the level. The method of
§\ref{sec:method} integrates both (\code{placements.json}), each ray at $h = \val{H}$
and at $2h$.
\begin{enumerate}[label=\arabic*.]
  \item $Q = \val{Qlo}$ at $(\val{xLo}, \val{yLo})$ only. The charge sits below the
    beam and pushes it upward. The $\val{angDown}^\circ$ ray and the axis ray leave
    through the right face of the world, above the detector. The $\val{angUp}^\circ$
    ray leaves through the top face of the world.

    The $\val{angDown}^\circ$ ray leaves at $(\val{LowDownX}, \val{LowDownY})$, at
    $t = \val{LowDownT}$.

    The momentum is $(\val{LowDownPx}, \val{LowDownPy})$, at $\val{LowDownAng}^\circ$
    to $+x$. The ordinate lies above $\val{yDhi}$. The estimate is $\val{LowDownErr}$.
    The radiated energy is $\val{LowDownRad}$.

    The axis ray leaves at $(\val{LowAxisX}, \val{LowAxisY})$, at
    $t = \val{LowAxisT}$.

    The momentum is $(\val{LowAxisPx}, \val{LowAxisPy})$, at $\val{LowAxisAng}^\circ$.
    The ordinate lies above $\val{yDhi}$. The estimate is $\val{LowAxisErr}$. The
    radiated energy is $\val{LowAxisRad}$.

    The $\val{angUp}^\circ$ ray leaves at $(\val{LowUpX}, \val{LowUpY})$, at
    $t = \val{LowUpT}$.

    The momentum is $(\val{LowUpPx}, \val{LowUpPy})$, at $\val{LowUpAng}^\circ$. The
    ordinate is the top of the world. The estimate is $\val{LowUpErr}$. The radiated
    energy is $\val{LowUpRad}$.
  \item $Q = \val{Qhi}$ at $(\val{xHi}, \val{yHi})$ only. The charge sits above the
    beam and pushes it downward. The $\val{angDown}^\circ$ ray and the axis ray leave
    through the right face of the world, below the detector. The $\val{angUp}^\circ$
    ray enters the detector outside the cone.

    The $\val{angDown}^\circ$ ray leaves at $(\val{HighDownX}, \val{HighDownY})$, at
    $t = \val{HighDownT}$.

    The momentum is $(\val{HighDownPx}, \val{HighDownPy})$, at
    $\val{HighDownAng}^\circ$. The ordinate lies below $\val{yDlo}$. The estimate is
    $\val{HighDownErr}$. The radiated energy is $\val{HighDownRad}$.

    The axis ray leaves at $(\val{HighAxisX}, \val{HighAxisY})$, at
    $t = \val{HighAxisT}$.

    The momentum is $(\val{HighAxisPx}, \val{HighAxisPy})$, at
    $\val{HighAxisAng}^\circ$. The ordinate lies below $\val{yDlo}$. The estimate is
    $\val{HighAxisErr}$. The radiated energy is $\val{HighAxisRad}$.

    The $\val{angUp}^\circ$ ray enters the \val{HighUpFace} face at
    $(\val{HighUpX}, \val{HighUpY})$, at $t = \val{HighUpT}$.

    The momentum is $(\val{HighUpPx}, \val{HighUpPy})$, at $\val{HighUpAng}^\circ$.
    That angle lies outside the cone of half-angle $\val{halfDeg}^\circ$, so the
    outcome is Rejected. The reported state is the entry. The estimate is
    $\val{HighUpErr}$. The radiated energy is $\val{HighUpRad}$.
\end{enumerate}

The reference needs both charges. The lower one turns the rays upward; alone, they
leave above the detector. The upper one turns them downward; alone, two leave below
the detector and the $\val{angUp}^\circ$ ray enters while still heading down, outside
the cone. Together the three rays of Figure~\ref{fig:board} meet the box and the cone.
Section~\ref{sec:comparison} records the outcomes.

\section{What the game computes}\label{sec:engine}

The game integrates \xeqref{foundations}{eq:eom} with its eighth-order Runge--Kutta
method at the tolerances $10^{-10}$ and $10^{-12}$ (\physref{5}). It finds the crossing
on its interpolant (\physref{6}) and applies \xeqref{injection}{eq:cone} to the
velocity at that entry (\physref{6.1}). It verifies the outcome by comparing its two
runs (\physref{7}). The model is this document's, including \ref{as:radiation}
(\physref{8}). The agreement in §\ref{sec:comparison} is between the fixed-step flight
of §\ref{sec:method}, the cone of §\ref{sec:cone}, and that integration.

\section{Further reading}\label{sec:reading}

\begin{itemize}
  \item \levelref{injection}: the cone \xeqref{injection}{eq:cone}, applied here to
    three rays.
  \item \levelref{einzel_lens}: the two-charge potential and the kinetic energy.
  \item \levelref{around_the_wall}: the force of several centres and the fixed-step
    method with the enlarged-box test, used here at $N = \val{nCharges}$.
  \item \levelref{slingshot}: why two fixed centres have no closed relativistic orbit,
    and the Hermite crossing.
  \item \levelref{first_bend}: the acceleration that enters the radiated power.
  \item \xsecref{foundations}{sec:board} and \xsecref{foundations}{sec:numerics}: the
    order of events, and the step-halving estimate.
\end{itemize}

\levelfileappendix

\end{document}
